\chapter{Software is not data --- and why ``beyond FAIR''}
\label{ch:beyond-fair}

Research software differs from data in kind: it is executable \emph{and} human-readable, it
evolves over time, and it sits atop a deep web of dependencies. The FAIR principles were
designed for data; this chapter shows where they fit research software and where they break,
and introduces the CODE-beyond-FAIR roadmap~\autocite{NatureSD2026}.

\takeaway{FAIR is a starting point, not the destination, for living software.}

\section{The living nature of software}
% TODO draft from: swh-ardc.org::#swnotdata; source-code-different-short.org.
% Figure: a real dependency graph (matplotlib) to make "deep web of dependencies" concrete.
% Assets: common/images/dependency-xkcd-2347.png (CC-BY-NC, attribute) or supply-chain-deps.png;
% better, generate a matplotlib dep graph for a CC-BY-clean figure.

\section{FAIR for software: fit and friction}
% TODO draft: the SNS "What about FAIR?" argument. Refs: barker2022fair4rs (FAIR4RS), 2020GtCitation.
% KEY POINT: the I (interoperable) and R (reusable) of FAIR already have precise, established
% meanings in software engineering --- and they differ from FAIR's data-centric intent. An
% experienced developer hears "interoperable/reusable" and thinks APIs, ABIs, modularity,
% portability --- not metadata. So the FAIR vocabulary actively CONFUSES practitioners. Keep
% what FAIR gets right (Findable + Accessible, delivered by archiving + intrinsic identifiers);
% be explicit that software interoperability/reusability is a distinct, harder, long-studied
% problem, not a metadata checkbox.
\takeaway{For software, \emph{interoperable} and \emph{reusable} already mean something precise
--- and it is not what FAIR means by them.}

\section{The CODE-beyond-FAIR roadmap}
% TODO: tiered roadmap for reusable research software. Ref: NatureSD2026.
